% RESEARCH_DRAFT_SHA256: 26e34d51305445d1551ed6dcaed50f83a14c2a9f80f413f352aef837bdbc2c8c
% Layout and source-reference wrappers expanded for the web renderer.
\documentclass[11pt]{article}
\usepackage[margin=25mm]{geometry}
\usepackage{fontspec,unicode-math}
\setmainfont{Latin Modern Roman}
\setmathfont{Latin Modern Math}
\usepackage{amsmath,amssymb,mathtools,xurl,hyperref}
\hypersetup{colorlinks=true,urlcolor=blue,linkcolor=blue}
\setlength{\parindent}{0pt}
\setlength{\parskip}{5pt}
\setlength{\emergencystretch}{3em}
\newcommand{\sref}[2]{\hyperlink{source:#1:#2}{[S#2]}}
\newcommand{\cataloguescope}[1]{\paragraph{Source question.}#1}
\newcommand{\floor}[1]{\left\lfloor#1\right\rfloor}
\begin{document}

\section*{Generated linear factors: a modulus-degree bound and failure of unrestricted uniformity}
\noindent catalogue entry; catalogue entry. Research attempt dated 11 October 2026.
\paragraph{Result in scope.}This manuscript contains a source-specific attempt and complete proofs of the stated auxiliary calculations. It makes no new solution or novelty claim for the entire target.
\section{Original problem and preserved definition}
The following definition and source list are reproduced from the repository. Any earlier inspection dates and reported literature claims in this preserved text describe that source record, not additional verification by this attempt.
% BEGIN PRESERVED DEFINITION
\subsection{Definitions and mathematical statement}

\paragraph{Shared notation.}Standard number systems are $\mathbb N,\mathbb Z,\mathbb Q,\mathbb R,\mathbb C$. Specialized notation and index ranges follow the cited source and the definitions below.


For a prime $p$, $\mathbb F_p[X]$ is the polynomial ring over the field of $p$ elements. In a quotient ring $\mathbb F_p[X]/(h)$, the residue class of $X-s$ is a unit when it is coprime to $h$. These units can generate a finite multiplicative subgroup whose order depends on the modulus and the chosen elements. Group-size lower bounds must state that ambient quotient and the degree and factorization assumptions on $h$.
\paragraph{Record scope and qualifications.}
The catalogue title is retained as metadata. Named certificates, contracts, bounds or frameworks in that title are not asserted as results of the original source extract.
The extraction drops cardinality symbols and does not explicitly give the ambient quotient or necessary modulus hypotheses. The numerical bounds in the historical passage are preserved as source claims rather than asserted for unrestricted $h$. The supplied catalogue marks this record partially solved; its present research status was not independently checked.
The source text below is rendered literally where necessary; damaged formulas and composite extracts are not silently repaired into a new mathematical claim.
\paragraph{Statement recorded in the supplied source.}
Question 2. Consider Fp and d ∈ Z>0, and S ⊂ Fp with S = d. Let h(X) ∈

Fp[X] with h(s) 6 = 0 for all s ∈ S. Let G be the group generated by X − s for s ∈ S. In the original AKS paper, we have G ≥ 2d. There are techniques for getting G ≥ (5.82) d

(using lattice point counting); there seems to be much more room for improvement (perhaps using ABC). (Bernstein, Voloch) Perhaps (log G)/d → ∞

as d → ∞ uniformly over h. \href{https://aimath.org/WWN/primesinp/primesinp.pdf}{[S2]}
\paragraph{Source status.}
The supplied catalogue labels this source record partially solved; that metadata was checked against the supplied snapshot on 2026-10-11. The label is preserved for provenance and is not an independent verification of present open status or mathematical correctness.
\subsection{Short English statement}
The source asks how large a group generated by linear factors can be in an AKS-related setting.
\subsection{Sources}
\begin{itemize}
\item[S1] Ulam AI and the UnsolvedMath Contributors, UnsolvedMath: A Curated Collection of Open Mathematics Problems, catalogue entry. Catalogue curation and metadata: CC BY 4.0. The source extract is retained; definitions, English summaries and qualifications are editorial additions.
\url{https://huggingface.co/datasets/ulamai/UnsolvedMath}.
\item[S2] AIM Workshop Problem Lists, source linked by the catalogue.
\url{https://aimath.org/WWN/primesinp/primesinp.pdf}.
\end{itemize}
% END PRESERVED DEFINITION
\subsection{Verbatim supplied catalogue question}
Question 2. Consider Fp and d ∈ Z>0, and S ⊂ Fp with S = d. Let h(X) ∈

Fp[X] with h(s) 6 = 0 for all s ∈ S. Let G be the group generated by X − s for s ∈ S. In the original AKS paper, we have G ≥ 2d. There are techniques for getting G ≥ (5.82) d

(using lattice point counting); there seems to be much more room for improvement (perhaps using ABC). (Bernstein, Voloch) Perhaps (log G)/d → ∞

as d → ∞ uniformly over h.

\section{Ambient convention and original scope}
The primary report's remarks identify the intended ambient ring as
$\mathbb F_p[X]/(h)$. Its question does not specify the degree or
factorization hypotheses on $h$ needed for the cited exponential
bounds. We make one literal finite-ring interpretation and check its
quantifiers. Let $h$ be nonconstant of degree $m$, let $S\subseteq
\mathbb F_p$ have size $d\ge1$, and suppose $h(s)\ne0$ for every
$s\in S$. Then $X-s$ is coprime to $h$: division by $X-s$ has
remainder $h(s)$. Bezout's identity in the polynomial Euclidean
algorithm makes each residue class a unit. Let $G$ be their generated
subgroup of this finite unit group.

\paragraph{Unrestricted uniform divergence is false in this reading.}
Choose any $a\in\mathbb F_p$, $h=X-a$, and
$S=\mathbb F_p\setminus\{a\}$. Evaluation at $a$ identifies the
quotient with $\mathbb F_p$, and the generators $a-s$ run through
all nonzero elements. Hence $G=\mathbb F_p^*$ and $|G|=p-1=d$.
Along arbitrarily large primes,
$\log|G|/d=\log d/d\to0$, rather than infinity.
There are arbitrarily large primes since the product of any finite
list plus one has a prime divisor outside the list. The limit follows,
for example, from $\log d\le\sqrt d$ for all sufficiently large $d$,
verified by differentiating $\sqrt d-\log d$ beyond four.
This refutes only the unrestricted degree-one-modulus version. It does
not refute an unstated AKS condition on $h$, and is not advertised as
a resolution of the intended strengthened group-size problem.

\paragraph{A complete degree-sensitive lower bound.}
All monic products
\[
P_e(X)=\prod_{s\in S}(X-s)^{e_s},\qquad
 e_s\ge0,\quad\sum_s e_s<m,
\]
have distinct residues modulo $h$. If two residues agreed, $h$ would
divide their difference, a polynomial of degree below $m$; hence the
difference would be zero. Unique factorization into the distinct
linear factors $X-s$ then forces every exponent to agree. The latter
can also be seen without a general factorization theorem by repeatedly
dividing out $X-s$ and evaluating at $s$ to recover its multiplicity.
Each product is a unit and lies in $G$. The number of exponent vectors
is
\[
|G|\ge\binom{d+m-1}{d}.
\]
To count them, introduce a slack exponent
$e_0=m-1-\sum_s e_s\ge0$. Weak compositions of $m-1$ into $d+1$
parts correspond to strings of $m-1$ identical stars and $d$ separators,
so their number is the displayed binomial coefficient.
In particular
\[
\binom{d+m-1}{d}=\prod_{j=0}^{d-1}\frac{m+j}{j+1}
 \ge(m/d)^d.
\]
Thus the desired divergence does follow in the additional regime
$m/d\to\infty$. This statement is conditional on that new regime;
it is not asserted uniformly over the original unspecified moduli.

\paragraph{Where the direct counting route stops.}
Once total degree reaches $m$, products can coincide modulo $h$ and
the proof no longer gives injectivity. The extreme $m=1$ example shows
that ignoring these collisions is genuinely invalid. In the difficult
regime $m$ comparable to $d$, the bound proves only exponential growth
with bounded base and therefore does not establish
$\log|G|/d\to\infty$. No assertion about the historical base $5.82$
is made without its additional hypotheses. A new way to control
higher-degree multiplicative collisions would be required.

\section{Audit and source status}
The primary Question 2 and its quotient-ring remarks were inspected.
The audit checked units, monic representatives, the modulus degree,
all asymptotic parameters, and the difference between a literal missing-
hypothesis counterexample and an intended AKS claim. No numerical
experiment or group computation was run. These elementary bounds and
the linear-modulus example are known-type scope diagnostics, with no
new resolution of the intended general question claimed.

\section{Process and contributions}
Mathematics and model-assisted adversarial self-audit: OpenAI gpt-6-astra, ultra reasoning, as dispatched for this research batch; exact serving revision not exposed. Codex (GPT-6; exact serving revision not exposed) materially assisted with orchestration and format review. No human mathematical authorship, mathematical edits, or independent human proof checking is asserted.
Contribution prompt A was applied to the full definition and supplied source row, then prompt B was applied adversarially to the derived results. The item-specific audit above records the scope checks and rejected extrapolations. Local assembly and structural checks do not establish mathematical correctness. The source catalogue status is preserved in TOP\_PROBLEM and is separate from this attempt status.
\paragraph{Integration review.} Codex (GPT-6; exact serving revision not exposed) reviewed the written arguments and expanded layout and source-reference wrappers. Repository metadata and the production web renderer were checked using Python 3.13 and JavaScript with MathJax 3.2.2. This model-assisted review is not independent human proof verification; no PDF compilation or formal proof checking is asserted.

\section{Completion Estimate}
\noindent\textbf{COMPLETION: 0\%}
The proved material is a rederivation of known or auxiliary facts rather than a substantiated new advance on the original target. The unrestricted reconstruction has a degree-one counterexample and a general degree-sensitive bound, but the intended AKS modulus hypotheses and difficult comparable-degree regime remain unresolved.
\section{References}
The complete original Sources subsection above is retained, including catalogue attribution and licence notices. The additional list records primary sources actually consulted for this attempt on 11 October 2026; uncited source links reproduced above were not all rechecked.
\begin{itemize}
\item American Institute of Mathematics, Future directions in algorithmic number theory, workshop 24--28 March 2003, report version 30 April 2003; original question, relevant lecture context and adjoining remarks inspected on 11 October 2026. \url{https://aimath.org/WWN/primesinp/primesinp.pdf}.
\end{itemize}
\section{Final status}
\textbf{Unresolved in this attempt.} The unrestricted reconstruction has a degree-one counterexample and a general degree-sensitive bound, but the intended AKS modulus hypotheses and difficult comparable-degree regime remain unresolved.
\end{document}
